%------------------------------------------------------------------------------%

\usepackage{apresentacao}
\usepackage{mathconfig}
\usepackage{tabto}

%\setlength{\parskip}{0.5\baselineskip}

\newcounter{exemplo}
\setcounter{exemplo}{0}

\makeatletter
\graphicspath  {{./figs/}}
\def\input@path{{./figs/}}
\makeatother

\newcommand{\D}[2]{\dfrac{\partial {#1}}{\partial {#2}}}

\newcommand{\vetor}[2]{\left(\begin{array}{c} #1 \\ #2 \end{array}\right)}

\newcommand{\veci}{\bm{i}}
\newcommand{\vecj}{\bm{j}}
\newcommand{\veck}{\bm{k}}

%------------------------------------------------------------------------------%
\title {Derivada Direcional}
\author{\large Luis Alberto D'Afonseca}
\date  {Cálculo de Funções de Várias Variáveis I}

%------------------------------------------------------------------------------%
\begin{document}

%------------------------------------------------------------------------------%
\begin{frame}
    \titlepage
\end{frame}

%------------------------------------------------------------------------------%
\section{Derivada Direcional}

%------------------------------------------------------------------------------%
\framedgraphic{Derivada Direcional}{derivada_direcional}{}

%------------------------------------------------------------------------------%
\begin{frame}{Definição}
  A derivada de $f$ em $(x_0, y_0)$
  \pause
  na direção do vetor unitário
  \[
    u
    \pause
    = u_1 \veci + u_2 \vecj
    \pause
    = \vetor{u_1}{u_2}
  \]
  \pause
  é o número
  \[
    \left(\frac{df}{ds}\right)_{u, (x_0, y_0)}
    \pause
    = \left(D_u f\right)_{(x_0, y_0)}
    \pause
    = \lim_{s\to 0} \frac{f(x_0+su_1, y_0+su_2)-f(x_0, y_0)}{s}
  \]
\end{frame}


%------------------------------------------------------------------------------%
\framedgraphic{Interpretação da Derivada Direcional}{derivada_direcional_interpretacao}{}

%------------------------------------------------------------------------------%
\begin{frame}{Interpretação da Derivada Direcional}
  A derivada direcional é a taxa de variação da função $f$ na direção do vetor $u$

  \pause
  \vspace{\baselineskip}
  Se \(u=\veci=\vetor{1}{0}\)
  \pause
  então
  \(
    \left(D_u f\right) = \D{f}{x}
  \)

  \pause
  \vspace{\baselineskip}
  Se \(u=\vecj=\vetor{0}{1}\)
  \pause
  então
  \(
    \left(D_u f\right) = \D{f}{y}
  \)
\end{frame}

%------------------------------------------------------------------------------%
\addtocounter{exemplo}{1}
\begin{frame}{Exemplo \theexemplo}
  Use a definição para calcular a derivada de \(f(x, y) = x^2 + xy\)
  no ponto \((1, 2)\) na direção
  \(
    u = 2\veci + 3\vecj
  \)

  \vspace{2\baselineskip}
  (Para simplificar os cálculos, ignore o fato desse vetor não ser unitário)
\end{frame}

\begin{frame}{Exemplo \theexemplo}
  \begin{align*}
    \onslide<+->{\left(\frac{df}{du}\right)_{u, (1, 2)}}
    \onslide<+->{& = \lim_{s\to 0} \frac{f(x_0+su_1, y_0+su_2)-f(x_0, y_0)}{s} \\}
    \onslide<+->{& = \lim_{s\to 0} \frac{f(1+2s, 2+3s)-f(1, 2)}{s} \\[2mm]}
    \onslide<+->{& = \lim_{s\to 0} \frac{\big((1+2s)^2 + (1+2s)(2+3s)\big)
                                         -\left(1^2+1\times 2\right)}{s} \\[2mm]}
    \onslide<+->{& = \lim_{s\to 0} \frac{(1 + 4s + 4s^2) + (2 + 3s + 4s + 6s^2) -3}{s} \\[2mm]}
    \onslide<+->{& = \lim_{s\to 0} \frac{11s + 10s^2}{s} \\[2mm]}
    \onslide<+->{& = \lim_{s\to 0} (11 + 10s)}
    \onslide<+->{  = 11}
  \end{align*}
\end{frame}

%------------------------------------------------------------------------------%
\section{Vetor Gradiente}

%------------------------------------------------------------------------------%
\begin{frame}{Vetor Gradiente}
  O \emph{vetor gradiente} de $f(x, y)$ em um ponto $(x_0, y_0)$
  \pause
  é
  \[
    \nabla f
    \pause
    \;=\; \D{f}{x}\veci \,+\, \D{f}{y}\vecj
    \pause
    \;=\; \left(\begin{array}{c} \D{f}{x} \\[5mm] \D{f}{y} \end{array}\right)
  \]
\end{frame}

%------------------------------------------------------------------------------%
\addtocounter{exemplo}{1}
\begin{frame}{Exemplo \theexemplo}
  Calcule o vetor gradiente da função \(f(x, y) = \sqrt{2x+3y}\)
  no ponto \((-1, 2)\)
\end{frame}

\begin{frame}{Exemplo \theexemplo}
  Calculando as derivadas parciais

  \pause
  \[
    \D{f}{x}
    \pause
    \;=\; \D{}{x}\sqrt{2x+3y}
    \pause
    \;=\; \D{}{x}\left(2x+3y\right)^{\sfrac{1}{2}}
  \]
  \[
    \hphantom{\D{f}{x}}
    \pause
    \;=\; \frac{1}{2}\left(2x+3y\right)^{\sfrac{-1}{2}}\D{}{x}\left(2x+3y\right)
    \pause
    \;=\; \frac{1}{\sqrt{2x+3y}}
  \]

  \pause
  \[
    \D{f}{x}(-1, 2)
    \pause
    \;=\; \frac{1}{\sqrt{2(-1)+3(2)}}
    \pause
    \;=\; \frac{1}{\sqrt{-2+6}}
    \pause
    \;=\; \frac{1}{\sqrt{4}}
    \pause
    \;=\; \frac{1}{2}
  \]
\end{frame}

\begin{frame}{Exemplo \theexemplo}
  Calculando as derivadas parciais

  \pause
  \[
    \D{f}{y}
    \pause
    \;=\; \D{}{y}\sqrt{2x+3y}
    \pause
    \;=\; \D{}{y}\left(2x+3y\right)^{\sfrac{1}{2}}
    \]
    \[
    \hphantom{\D{f}{y}}
    \pause
    \;=\; \frac{1}{2}\left(2x+3y\right)^{\sfrac{-1}{2}}\D{}{y}\left(2x+3y\right)
    \pause
    \;=\; \frac{3}{2\sqrt{2x+3y}}
  \]

  \pause
  \[
    \D{f}{y}(-1, 2)
    \pause
    \;=\; \frac{3}{2\sqrt{2(-1)+3(2)}}
    \pause
    \;=\; \frac{3}{2\sqrt{-2+6}}
    \pause
    \;=\; \frac{3}{2\sqrt{4}}
    \pause
    \;=\; \frac{3}{4}
  \]
\end{frame}

\begin{frame}{Exemplo \theexemplo}
  O vetor gradiente
  \[
    \nabla f
    \pause
    \;=\; \left(\begin{array}{c} \D{f}{x} \\[5mm] \D{f}{y} \end{array}\right)
    \pause
    \;=\; \left(\begin{array}{c}
             \dfrac{1}{\sqrt{2x+3y}} \\[5mm]
             \dfrac{3}{2\sqrt{2x+3y}}
          \end{array}\right)
    \pause
    \;=\; \frac{1}{\sqrt{2x+3y}} \left(\begin{array}{c}
            1 \\[2mm]
            \sfrac{3}{2}
          \end{array}\right)
  \]
\end{frame}

\begin{frame}{Exemplo \theexemplo}
  O vetor gradiente no ponto \((-1, 2)\)
  \[
    \nabla f(-1, 2)
    \;=\; \left[\frac{1}{\sqrt{2x+3y}} \left(\begin{array}{c}
      1 \\[2mm]
      \sfrac{3}{2}
    \end{array}\right)\right]\evalat{(-1, 2)}{}
    \;=\; \vetor{\sfrac{1}{2}}{\sfrac{3}{4}}
  \]
\end{frame}

%------------------------------------------------------------------------------%
\section{Calculando a Derivada Direcional}

%------------------------------------------------------------------------------%
\begin{frame}{Calculando a Derivada Direcional}
  Considere a reta que passa por \((x_0, y_0)\)
  \pause
  na direção de \(u=u_1\veci+u_2\vecj\)
  \pause
  \[
    x = x_0 + su_1 \qquad
    y = y_0 + su_2
  \]

  Restringindo $f$ à reta
  \[
    f(s) = f\big(x(s), y(s)\big)
  \]

  \pause
  Pela regra da cadeira temos
  \[
    D_u f
    \pause
    \;=\; \frac{df}{ds}
    \pause
    \;=\; \D{f}{x}\frac{dx}{ds} + \D{f}{y}\frac{dy}{ds}
    \pause
    \;=\; \D{f}{x}u_1 + \D{f}{y}u_2
    \pause
    \;=\; \left(\begin{array}{c} \D{f}{x} \\[5mm] \D{f}{y} \end{array}\right)
      \cdot \vetor{u_1}{u_2}
    \pause
    \;=\; \nabla f \cdot u
  \]
\end{frame}

%------------------------------------------------------------------------------%
\addtocounter{exemplo}{1}
\begin{frame}{Exemplo \theexemplo}
  Encontre a derivada de \(f(x, y) = xe^y + \cos(xy)\)
  no ponto \((2, 0)\)
  na direção \(v=3\veci-4\vecj\)
\end{frame}

\begin{frame}{Exemplo \theexemplo}
  O vetor $v$ não é unitário,
  \pause
  então precisamos dividi-lo pelo seu
  comprimento
  \[
    \pause
    u
    \;=\; \frac{v}{\abs{v}}
    \pause
    \;=\; \frac{3\veci-4\vecj}{\sqrt{3^2+4^2}}
    \pause
    \;=\; \frac{3\veci-4\vecj}{\sqrt{25}}
    \pause
    \;=\; \frac{3}{5}\veci - \frac{4}{5}\vecj
  \]
\end{frame}

\begin{frame}{Exemplo \theexemplo}
  Calculando as derivadas parciais

  \pause
  \[
    \D{f}{x}
    \pause
    \;=\;  \D{}{x}\left(xe^y + \cos(xy)\right)
    \pause
    \;=\;  e^y - y\sin(xy)
  \]

  \pause
  \[
    \D{f}{x}(2, 0)
    \pause
    \;=\;  e^0 - 0\sin(2\times 0)
    \pause
    \;=\;  1
  \]
\end{frame}

\begin{frame}{Exemplo \theexemplo}
  Calculando as derivadas parciais

  \pause
  \[
    \D{f}{y}
    \pause
    \;=\;  \D{}{x}\left(xe^y + \cos(xy)\right)
    \pause
    \;=\;  xe^y - x\sin(xy)
  \]

  \pause
  \[
    \D{f}{y}(2, 0)
    \pause
    \;=\;  2e^0 - 2\sin(2\times 0)
    \pause
    \;=\;  2
  \]
\end{frame}

\begin{frame}{Exemplo \theexemplo}
  Calculando a derivada direcional
  \[
    \left(D_u f\right)\evalat{(2, 0)}{}
    \pause
    \;=\; \left(\nabla f\right)\evalat{(2, 0)}{} \cdot u
    \pause
    \;=\; \vetor{1}{2} \cdot \vetor{\sfrac{3}{5}}{-\sfrac{4}{5}}
  \]
  \[
    \hphantom{\left(D_u f\right)\evalat{(2, 0)}{}}
    \pause
    \;=\; 1 \frac{3}{5} + 2 \frac{-4}{5}
    \pause
    \;=\; \frac{3}{5} - \frac{8}{5}
    \pause
    \;=\; - \frac{5}{5}
    \pause
    \;=\; -1
  \]
\end{frame}

%------------------------------------------------------------------------------%
\section{Lista Mínima}

%------------------------------------------------------------------------------%
\begin{frame}{Lista Mínima}
  Cálculo Vol. 2 do Thomas 12\textsuperscript{a} ed. -- Seção 14.5
  \begin{enumerate}
    \item Estudar o texto da seção
  \end{enumerate}
  \vfill
  Atenção: A prova é baseada no livro, não nas apresentações
\end{frame}

%------------------------------------------------------------------------------%
\begin{frame}
  \tableofcontents
\end{frame}

%------------------------------------------------------------------------------%
\end{document}
%------------------------------------------------------------------------------%